\documentclass[10pt,a4paper]{amsart}
\usepackage{amsmath,amssymb,amsthm,graphicx}
% Retained minimal preamble for the Stacks Project chapter flat.tex (More on Flatness).
% Verbatim excerpt of the chapter preamble: Xy-pic setup (native preamble lines 14-18) and the native
% theorem declarations with the mathematical macros (native preamble lines 154-217). Copied mechanically;
% nothing added, removed or redefined. The wrapper class and the amsmath/amssymb/amsthm/graphicx base
% packages are supplied by the assembler, so the class/externaldocument part of the native preamble is omitted.
% For commutative diagrams we use Xy-pic
\usepackage[all]{xy}

% We use 2cell for 2-commutative diagrams.
\xyoption{2cell}
% Theorem environments.
%
\theoremstyle{plain}
\newtheorem{theorem}[subsection]{Theorem}
\newtheorem{proposition}[subsection]{Proposition}
\newtheorem{lemma}[subsection]{Lemma}

\theoremstyle{definition}
\newtheorem{definition}[subsection]{Definition}
\newtheorem{example}[subsection]{Example}
\newtheorem{exercise}[subsection]{Exercise}
\newtheorem{situation}[subsection]{Situation}

\theoremstyle{remark}
\newtheorem{remark}[subsection]{Remark}
\newtheorem{remarks}[subsection]{Remarks}

\numberwithin{equation}{subsection}

% Macros
%
\def\lim{\mathop{\mathrm{lim}}\nolimits}
\def\colim{\mathop{\mathrm{colim}}\nolimits}
\def\Spec{\mathop{\mathrm{Spec}}}
\def\Hom{\mathop{\mathrm{Hom}}\nolimits}
\def\Ext{\mathop{\mathrm{Ext}}\nolimits}
\def\SheafHom{\mathop{\mathcal{H}\!\mathit{om}}\nolimits}
\def\SheafExt{\mathop{\mathcal{E}\!\mathit{xt}}\nolimits}
\def\Sch{\mathit{Sch}}
\def\Mor{\mathop{\mathrm{Mor}}\nolimits}
\def\Ob{\mathop{\mathrm{Ob}}\nolimits}
\def\Sh{\mathop{\mathit{Sh}}\nolimits}
\def\NL{\mathop{N\!L}\nolimits}
\def\CH{\mathop{\mathrm{CH}}\nolimits}
\def\proetale{{pro\text{-}\acute{e}tale}}
\def\etale{{\acute{e}tale}}
\def\QCoh{\mathit{QCoh}}
\def\Ker{\mathop{\mathrm{Ker}}}
\def\Im{\mathop{\mathrm{Im}}}
\def\Coker{\mathop{\mathrm{Coker}}}
\def\Coim{\mathop{\mathrm{Coim}}}

% Boxtimes
%
\DeclareMathSymbol{\boxtimes}{\mathbin}{AMSa}{"02}

%
% Macros for moduli stacks/spaces
%
\def\QCohstack{\mathcal{QC}\!\mathit{oh}}
\def\Cohstack{\mathcal{C}\!\mathit{oh}}
\def\Spacesstack{\mathcal{S}\!\mathit{paces}}
\def\Quotfunctor{\mathrm{Quot}}
\def\Hilbfunctor{\mathrm{Hilb}}
\def\Curvesstack{\mathcal{C}\!\mathit{urves}}
\def\Polarizedstack{\mathcal{P}\!\mathit{olarized}}
\def\Complexesstack{\mathcal{C}\!\mathit{omplexes}}
% \Pic is the operator that assigns to X its picard group, usage \Pic(X)
% \Picardstack_{X/B} denotes the Picard stack of X over B
% \Picardfunctor_{X/B} denotes the Picard functor of X over B
\def\Pic{\mathop{\mathrm{Pic}}\nolimits}
\def\Picardstack{\mathcal{P}\!\mathit{ic}}
\def\Picardfunctor{\mathrm{Pic}}
\def\Deformationcategory{\mathcal{D}\!\mathit{ef}}
\usepackage{cleveref}
\crefname{originallemma}{original lemma}{original lemmas}
\Crefname{originallemma}{Original lemma}{Original lemmas}
\makeatletter
\newlabel{algebra-lemma-ass-quotient-ring}{{05BY}{source}{}{sourcepointer-1}{}}
\newlabel{algebra-lemma-ass-quotient-ring@cref}{{[originallemma][0][]05BY}{[1][0][]source}{}{}{}}
\makeatother
\crefname{originallemma}{original lemma}{original lemmas}
\Crefname{originallemma}{Original lemma}{Original lemmas}
\makeatletter
\newlabel{algebra-lemma-finite-ass}{{00LC}{source}{}{sourcepointer-2}{}}
\newlabel{algebra-lemma-finite-ass@cref}{{[originallemma][0][]00LC}{[1][0][]source}{}{}{}}
\makeatother
\crefname{originallemma}{original lemma}{original lemmas}
\Crefname{originallemma}{Original lemma}{Original lemmas}
\makeatletter
\newlabel{algebra-lemma-ass-zero-divisors}{{00LD}{source}{}{sourcepointer-3}{}}
\newlabel{algebra-lemma-ass-zero-divisors@cref}{{[originallemma][0][]00LD}{[1][0][]source}{}{}{}}
\makeatother
\crefname{originallemma}{original lemma}{original lemmas}
\Crefname{originallemma}{Original lemma}{Original lemmas}
\makeatletter
\newlabel{algebra-lemma-spec-localization}{{00E3}{source}{}{sourcepointer-4}{}}
\newlabel{algebra-lemma-spec-localization@cref}{{[originallemma][0][]00E3}{[1][0][]source}{}{}{}}
\makeatother
\crefname{originallemma}{original lemma}{original lemmas}
\Crefname{originallemma}{Original lemma}{Original lemmas}
\makeatletter
\newlabel{algebra-lemma-silly}{{00DS}{source}{}{sourcepointer-5}{}}
\newlabel{algebra-lemma-silly@cref}{{[originallemma][0][]00DS}{[1][0][]source}{}{}{}}
\makeatother
\title{Localizing away from the zerodivisors of a fibre module: the spectrum is semi-local on the associated primes}
\author{The Stacks Project Authors}
\date{stacks-project:flat.tex:local-snapshot-20261011; GNU FDL 1.2 or later}
\begin{document}
\maketitle
\noindent\textit{Editorial note.} All delivered mathematical text is the author's own, copied from the pinned Stacks Project chapter file flat.tex; no statement or proof step was rewritten, reordered or removed.\par\medskip
\noindent\textit{Editorial note.} The preamble is the chapter's own retained minimal preamble (sources/stacks-flat/retained-preamble.tex), a verbatim excerpt of the native Xy-pic setup and of the native theorem declarations and mathematical macros; no macro was invented or redefined.\par\medskip
\noindent\textit{Editorial note.} Every referenced result of this card lives in an earlier chapter, so each locator points at the verbatim line of the project's own tag list staged in the snapshot (sources/stacks-chow/tags.txt) that carries the canonical tag.\par\medskip
\noindent\textit{Allowed standard prerequisites (editorial).}\begin{itemize}
\item Associated primes of a module over a quotient ring. Algebra, Lemma labelled algebra-lemma-ass-quotient-ring (canonical Stacks tag 05BY); invoked at proof line 1803. Supplies the identification Ass\_S(Nbar) = Ass\_Sbar(Nbar) for the module Nbar = N/m\_R N over Sbar = S/m\_R S; the rest of the argument is retained in the counted span.
\item Finiteness of the set of associated primes of a finite module over a Noetherian ring. Algebra, Lemma labelled algebra-lemma-finite-ass (canonical Stacks tag 00LC); invoked at proof lines 1804-1806. Supplies only the finiteness of Ass\_S(Nbar), which is what makes the enumeration q\_1, ..., q\_r available.
\item Zerodivisors are the union of the associated primes. Algebra, Lemma labelled algebra-lemma-ass-zero-divisors (canonical Stacks tag 00LD); invoked at proof lines 1807-1808. Supplies the identity Sigma = S minus the union of the q\_i, which is the translation the counted proof needs in order to describe the localisation.
\item Primes of a localisation and primes contained in a given prime. Algebra, Lemma algebra-lemma-spec-localization (tag 00E3) and algebra-lemma-silly (tag 00DS); invoked at proof lines 1809-1813. Supply the description of Spec(Sigma\textasciicircum{}\{-1\}S) and the behaviour of primes under localisation; the resulting correspondence with the q\_i and the maximality argument are inside the counted span.
\end{itemize}\medskip
\ifdefined\hypertarget\hypertarget{sourcepointer-1}{}\fi
\noindent\textit{Source reference.} Original lemma 05BY: sources/stacks-chow/tags.txt, source lines 6901--6901. Cross-chapter standard prerequisite invoked once at proof line 1803 to identify associated primes over S with associated primes over S/m\_R S; it is a citation of a general algebra fact and replaces nothing in the retained proof.\par\medskip
\ifdefined\hypertarget\hypertarget{sourcepointer-2}{}\fi
\noindent\textit{Source reference.} Original lemma 00LC: sources/stacks-chow/tags.txt, source lines 782--782. Cross-chapter standard prerequisite invoked once at proof line 1805 for finiteness of the set of associated primes of a finite module over a Noetherian ring.\par\medskip
\ifdefined\hypertarget\hypertarget{sourcepointer-3}{}\fi
\noindent\textit{Source reference.} Original lemma 00LD: sources/stacks-chow/tags.txt, source lines 783--783. Cross-chapter standard prerequisite invoked once at proof line 1808 for the description of the zerodivisors of Nbar as the union of its associated primes.\par\medskip
\ifdefined\hypertarget\hypertarget{sourcepointer-4}{}\fi
\noindent\textit{Source reference.} Original lemma 00E3: sources/stacks-chow/tags.txt, source lines 520--520. Cross-chapter standard prerequisite invoked once at proof lines 1809-1810 for the description of the primes of a localisation; the translation into the primes contained in one of the q\_i is the counted step.\par\medskip
\ifdefined\hypertarget\hypertarget{sourcepointer-5}{}\fi
\noindent\textit{Source reference.} Original lemma 00DS: sources/stacks-chow/tags.txt, source lines 508--508. Cross-chapter standard prerequisite invoked once at proof lines 1811-1812 for comparing primes under localisation; the maximality conclusion for the maximal ideals of Sigma\textasciicircum{}\{-1\}S is retained in the counted span.\par\medskip
\begin{lemma}
\label{lemma-homothety-spectrum}
Let $R \to S$ be a ring map.
Let $N$ be a $S$-module.
Assume
\begin{enumerate}
\item $R$ is a local ring with maximal ideal $\mathfrak m$,
\item $\overline{S} = S/\mathfrak m S$ is Noetherian, and
\item $\overline{N} = N/\mathfrak m_R N$ is a finite $\overline{S}$-module.
\end{enumerate}
Let $\Sigma \subset S$ be the multiplicative subset of elements which are not
a zerodivisor on $\overline{N}$. Then $\Sigma^{-1}S$ is a semi-local ring
whose spectrum consists of primes $\mathfrak q \subset S$ contained in an
element of $\text{Ass}_S(\overline{N})$. Moreover, any maximal
ideal of $\Sigma^{-1}S$ corresponds to an associated prime of
$\overline{N}$ over $\overline{S}$.
\end{lemma}

\begin{proof}
Note that
$\text{Ass}_S(\overline{N}) = \text{Ass}_{\overline{S}}(\overline{N})$, see
Algebra, Lemma \ref{algebra-lemma-ass-quotient-ring}.
This is a finite set by
Algebra, Lemma \ref{algebra-lemma-finite-ass}.
Say $\{\mathfrak q_1, \ldots, \mathfrak q_r\} = \text{Ass}_S(\overline{N})$.
We have $\Sigma = S \setminus (\bigcup \mathfrak q_i)$ by
Algebra, Lemma \ref{algebra-lemma-ass-zero-divisors}.
By the description of $\Spec(\Sigma^{-1}S)$ in
Algebra, Lemma \ref{algebra-lemma-spec-localization}
and by
Algebra, Lemma \ref{algebra-lemma-silly}
we see that the primes of $\Sigma^{-1}S$ correspond to the primes of
$S$ contained in one of the $\mathfrak q_i$.
Hence the maximal ideals of $\Sigma^{-1}S$ correspond one-to-one with the
maximal (w.r.t.\ inclusion) elements of the set
$\{\mathfrak q_1, \ldots, \mathfrak q_r\}$. This proves the lemma.
\end{proof}

\end{document}
